%%%%%%%%%%%%
%
% $Autor: Wings $
% $Datum: 2019-03-05 08:03:15Z $
% $Pfad: ListOfParts.tex $
% $Version: 4250 $
% !TeX spellcheck = en_GB/de_DE
% !TeX encoding = utf8
% !TeX root = manual 
% !TeX TXS-program:bibliography = txs:///biber
%
%%%%%%%%%%%%

\chapter{Schrittweise Montage}
\subsubsection{Vorbereitung}
Entnehmen Die die vormontierten Materialien aus der Verpackung und überprüfen Sie die Unversehrtheit alles Teile. Wählen Sie einen geeigneten Ort für den Aufbau.
\subsubsection{Montage}
Stellen Sie den Blumentopf (Pos. 1) auf eine geeignete, ebene Fläche. Um Kratzer oder Beschädigungen an den Materialien zu vermeiden, empfiehlt es sich, eine weiche Unterlage wie z. B. eine Decke unter zulegen. Führen Sie anschließend die vormontierten Kabel, die sich am Deckel befinden, durch die dafür vorgesehene Kabeldurchführung im Blumentopf. Nun kann der Deckel mit dem bereits montierten Rohr und Sensor auf den Blumentopf aufgesetzt werden. Der Deckel ist passgenau gefertigt und rastet direkt auf dem Blumentopf ein, ohne dass eine zusätzliche Befestigung notwendig ist – dies ermöglicht eine einfache Abnahme bei Wartung oder Reinigung.

Nachdem der Deckel montiert ist, verbinden Sie das durchgeführte Kabel mit dem mitgelieferten E-Gehäuse. Stecken Sie dazu den Stecker durch die vorgesehene Öffnung in das Gehäuse. Der Sensor sowie die Pumpe sind nun elektrisch mit dem System verbunden.

Bevor das System an das Stromnetz angeschlossen wird, füllen Sie bitte etwa 5 Liter Wasser in den Behälter ein. Achten Sie darauf, dass der Wasserstand ausreichend ist, um die Pumpe vollständig zu bedecken. Im letzten Schritt können die beiliegenden Dekosteine gleichmäßig auf dem Deckel verteilt werden – diese dienen ausschließlich dekorativen Zwecken.

Sobald alles korrekt montiert und mit Wasser befüllt ist, schließen Sie das System an eine Steckdose an. Der Magic Faucet Fountain ist damit betriebsbereit.

